
\section{Multinational Enterprises}

\begin{frame}{Corporations, Organisations}
    \mytext{
    \begin{itemize}
        \item MNE multinational enterprise
        \item organisational structure
        \item key indicators: balance sheet, turnovers etc.
        \item markets served: how large is X as part of total
        \item headquarters, international presence, employees
        \item trends: past, present, future; announcements
        \item publications, lobbying
        \item main partners
        \item most important legislations
    \end{itemize}
    }
    \mycomment{
    While a company  website is a prime source, you should try to find additional, more critical information. 

    Do not fall for the advertisement. Look at the entire industry, and in the light of this look at the individual company.

    Some questions are: where is the company listed? What is the legal form, and who are the owners? In which currency does it publish its accounts? 
    
    Who are the main partners? On which prices and technologies will the fortune of the company or industry depend?
    }
\end{frame}



\begin{frame}{Market Forms}
    \mytext{
    \mytable{
    \renewcommand{\arraystretch}{2}
        \begin{tabular}{l||l|l|l|}
                    &   one seller  & few sellers   & many sellers \\\hline\hline                
        one buyer   &  bilat. mon.  & oligopsony     & monopsony     \\\hline
        many buyers &   monopoly    & oligopoly     & competition \\\hline
        \end{tabular}
        }
        
    % Note that "few" is relative:
    
    \begin{quotation}
        Four is few and six is many \cite{selten1994}
    \end{quotation}
    }
\end{frame}

\begin{frame}{Competition}
    \mytext{
    \begin{itemize}
        \item \textit{competere }lat. \textit{to strive together}
        \item \textit{concurrere} lat. \textit{to run together}
        \item not a zero-sum game
        \item think of a race (rather than combat)
        \item market entry barriers
        \item few participants: more combative 
        \item many participants: much intelligence
        \item property structures
    \end{itemize}
    }
    \mycomment{
    Not only does competition result in the largest possible quantities and the lowest possible price, it also incentivises sellers to strive for productivity increase, that is: to lower their costs. 

    However, perfect competition, as assumed in many models, is rare. There are typically barriers to entry.
    }
\end{frame}
 

\begin{frame}{Monopolies}
    \mytext{
    \begin{itemize}
        \item natural monopolies (e.g. networks)
        \item setting up monopolies (illegal in Europe)
        \item having a monopoly (illegal in the US)
        \item monopoly world wide or only at home?
        \item market entry?
        \item how to determine a monopoly?
        \item advantages of a monopoly?
        \item state monopoly, state run monopolies
        \item regulation, antitrust enforcement
    \end{itemize}    
    }
    \mycomment{
    Examples for natural monopolies are: port authorities, railway or energy system operators. However, questions remain. In particular, the monopolist can extend their market power to adjacent markets (e.g. Edison also selling light bulbs and only accepting their own light bulbs to be connected). 
    }
\end{frame}



\begin{frame}{Concentration Ratio}
    \mytext{
    The Concentration Ratio simply expresses the market share of the top $n$ companies:
    \begin{equation}
        CR_n := \sum_i^n a_i
    \end{equation}

    Scientifically better is the Herfindahl-Hirschman-Index, the sum of the squares of the market share of every firm:
    \begin{equation}
        HHI := \sum_{i=1}^N a_i^2
    \end{equation}
    with $a_i$ market share of participant $i$.
    HHI of 0 is perfect competition, 10000 is perfect monopoly. For technical reasons, often $N:=50$.
    }
\end{frame}




\begin{frame}{Predatory Pricing}
    \mytext{
        \begin{itemize}
            \item undercutting so that competitor can't survive
            \item barriers to entry: change is permanent
            \item monopoly rent afterwards allows to recoup losses
            \item often seen after new entry
            \item mostly illegal
            \item question of available capital
            \item state subsidies / global trade?
        \end{itemize}
    }
    \mycomment{
        Predatory pricing shifts the concurrent competition to a direct, interactive `game play': rather than to truly concentrate on one's own business and the customer, the actions of the \textbf{incumbent} or the \textbf{entrant} focuses on the effect on the other competitor(s).
    }
\end{frame}

\begin{frame}{State Monopolies}
    \mytext{
        \begin{itemize}
            \item State may be able to do what others can't
            \item Infrastructure
            \item expropriation
            \item drawbacks: under-investment
            \item privatisation
            \item re-communalisation
            \item socialism, monarchy...
        \end{itemize}
    }
    \mycomment{
        Here is room for debate. From an economic point of view, intelligent regulation will always beat communalisation, because the real incentive structure for all, including managers and / or politicians, must be taken into account.
    }
\end{frame}



\begin{frame}{Multinational Enterprises (MNEs)}
    \mytext{
    \begin{itemize}
        \item internationally invested
        \item trans-border transactions
        \item typically dependencies abroad
        \item concerned with currency, tariffs etc.
        \item upstream, downstream
        \item logistics, trade
        \item finance, IT
        \item ...
    \end{itemize}
    }
\end{frame}

\begin{frame}{Multinational Enterprises (MNEs)}
    \mytext{
    \begin{itemize}
        \item Dutch and British East India Companies
        \item Seven Sisters (oil): Shell, BP, Total...
        \item Unilever, Nestlé...
        \item JPMorgan Chase, HSBC, other Banks
        \item Keiretsu (Mitsubishi, ...)
        \item Magnificent Seven, ASML
        \item Volkswagen, Toyota, Stellantis...
        \item Bayer, BASF, Sinopec
        \item Walmart, Amazon
        \item Glencore
        \item Boeing, Airbus, Lockheed Martin
    \end{itemize}
    }
    \mycomment{
    In most markets, some multinational enterprises dominate the market somehow. Who are they - and what is their market share? What is the market share of the top 10 or top 50 companies? How concentrated is the market?
    }
\end{frame}


\begin{frame}{Company Analysis}
    \mycomment{
    You should approach the company analysis from a market perspective. Is it a bull market or a bear market? What is the share of the company? Is it a price taker? etc. 

    For the business perspective, check what the company has published. Publicly listed US companies are obliged to publish Annual Reports to Shareholders (ARS), or Form 10-K. In any legislation, you will find financial statements published---have a look at these.
    }
    \mytext{
    \begin{itemize}
        \item Financial Statements
        \item Balance sheets, income statement, cash flow
        \item Ownership structure, governance
        \item Assets and Liabilities
        \item Market share, turnover
        \item Cost structure, trends, policy
        \item Profitability, liquidity, leverage
        \item Comparable Company Analysis: benchmarking
        \item Technological Analysis, Life Cycle Analysis
    \end{itemize}
    }
\end{frame}


\begin{frame}{Organisation Structure}
    \showmytikz{../shared/tikz/tikz-owner-holding-company}
    \mycomment{
    What is the ownership (structure) of the company? Are there majority share holders? Is it part of a larger group or conglomerate?  Are there daughter companies? Are states invested?
    }
\end{frame}

\begin{frame}{Organisation Structure}
    \showmytikz{../shared/tikz/tikz-line-and-staff}
    \mycomment{
    How is the company organised? Which divisions and branches exist?

    Also: what is the ownership of the company? Is it part of a larger conglomerate? Is it part of a holding?

    Can you figure out how the company is structured: is it a traditional, hierarchical structure, is it project based, flat hierarchies, is it a matrix organisation?
    }
\end{frame}

\begin{frame}{Porter: Value Chain}
    \mycomment{In our context, we are not so much interested in the entire company, as complex as that structure may be, but more about the value chain within the market the company is covering. Is the company operating upstream or downstream? is it processing goods it owns, or is the processing a service? etc.}
    \showmytikz{../shared/tikz/tikz-porter-valuechain.tex}    
\end{frame}

\begin{frame}{Porter: Six Forces}
    \mycomment{The Six Forces scheme can be used to analyse markets. However, there is more about markets that this can cover.

    Porters Six Forces (or sometimes: Five Forces) are a much over-used, but nevertheless useful tool to look at a company, or a branch of said company (a product).  But it should not be followed blindly.
    }
    \showmytikz{../shared/tikz/tikz-porter-sixforces}
\end{frame}

\begin{frame}{SWOT Analysis}
    \mycomment{SWOT Analysis is another common tool to look at branches of a company. In our context, it may not be the most useful tool.}
    \showmytikz{../shared/tikz/tikz-swot.tex}
\end{frame}

\begin{frame}{BGC Matrix}
    \mycomment{The BGC matrix is often used to compare different fields, activities or products / services of a company. It can also be used to compare a company to another.}
    \showmytikz{../shared/tikz/tikz-bgc-matrix}
\end{frame}

\begin{frame}{Product Life Cycle}
    \mycomment{
    The product life cycle is particularly important in markets with large technological changes, or where fashion changes fast. It is much less relevant for commodities, albeit these may follow the life cycles of products for which the commodity is used (who buys whale fat today?).

    There are relevant phases prior to market introduction (here: left of point $A$). Does the company invest much in research and development?
    }
    \showmytikz{../shared/tikz/tikz-productlifecycle}
\end{frame}

\begin{frame}{Expenditure over Time}
    \showmytikz{../shared/tikz/tikz-expenditureovertime}
    \mycomment{
    Much of companies' investment relies on capital investment, which occurs prior to marketing of products. What assets can you make out? Does the company own e.g. mining fields, drilling platforms, farms, factories?  
    }
\end{frame}

\begin{frame}{Upstream and Downstream Integration}
    \showmytikz{../shared/tikz/tikz-valuechain}
    \mycomment{
        Is the company invested in upstream, downstream, or both? 
        
        If it serves multiple stages of the value chain, the company is said to be vertically integrated. If it serves a variety of related products on the same level of the value chain, it is said to be horizontally integrated.
    }
\end{frame}

% \begin{frame}{Risk Register}
%     \mycomment{A risk register is a valuable tool to identify risks (and opportunities!) for any enterpriese.}
%     \showmytikz{../shared/tikz/tikz-riskregister}    
% \end{frame}


\begin{frame}{International Organisations}
    \mytext{

    Look at large organisations that serve a number of (member) enterprises:
    
    \begin{itemize}
        \item industry associations on national and international level
        \item governmental industry associations
        \item[] (chambers of commerce)
        \item inter-governmental organisations, such as within the UN system, OECD etc.
    \end{itemize}
    }
    \mycomment{
    Due to the topic of this class, we will not be looking at civil organisations or NGOs. However, technically speaking, industry associations are NGOs (non-governmental). 

    Large MNEs are also international organisations. However, international organisations such as trade and industry associations are typically not directly involved in the business or industry they supervise or support.
    }
\end{frame}


